% !TEX TS-program = pdflatex
% !TeX program = pdflatex
% !TEX encoding = UTF-8
% !TeX spellcheck = en_US

\documentclass[11pt, a4paper]{article}
%\usepackage{fullpage}
\usepackage[left=1cm,right=1cm,top=1cm,bottom=2cm]{geometry}
\usepackage[fleqn]{amsmath}
\usepackage{amssymb}
%\usepackage{indentfirst}
\usepackage[T1]{fontenc}
\usepackage[utf8]{inputenc}
\usepackage[french,english]{babel}
\usepackage{txfonts} 
\usepackage[]{graphicx}
\usepackage{multirow}
\usepackage{hyperref}
\usepackage{parskip}
\usepackage{multicol}
\usepackage{wrapfig}
\usepackage{amssymb}
\usepackage{longtable}
\usepackage{wasysym}


\usepackage{turnstile}%Induction symbole

\usepackage{tikz}
\usetikzlibrary{arrows, automata}
\usetikzlibrary{decorations.pathmorphing}

\renewcommand{\baselinestretch}{1}

\setlength{\parindent}{24pt}


\begin{document}

%\selectlanguage {french}
%\pagestyle{empty} 

\noindent
\begin{tabular}{ll}
\multirow{3}{*}{\includegraphics[width=1.5cm]{../extra/logo/esi.nlp.pdf}} & 
\'Ecole national Supérieure d'Informatique\\
& 2\textsuperscript{nd} year second cycle (2CSSID)\\
& NLP: Natural Language Processing (2022-2023)
\end{tabular}\\[.25cm]
\noindent\rule{\textwidth}{2pt}\\[-0.5cm]
\begin{center}
{\LARGE \textbf{Tutorial 01}}
\begin{flushright}
	Abdelkrime Aries
\end{flushright}
\end{center}\vspace{-0.5cm}
\noindent\rule{\textwidth}{2pt}

\section*{1. Language models}


\subsection*{1.1.}

Select right propositions :

\begin{longtable}{|p{.95\textwidth}|}
	\hline 
	\Square\ Some smoothing methods can fix out of vocabulary problem. \\
	\Square\ Given a unigram model trained on standard English, the perplexity of the sentence ``\textbf{Karim teaches a course}" is higher than that of ``\textbf{a course Karim teaches}". \\
	\Square\ Given a unigram model trained on standard English, the perplexity of the sentence ``\textbf{Karim teaches a course}" is equal to that of ``\textbf{a course Karim teaches}".\\
	\Square\ Given a unigram model trained on standard English, the perplexity of the sentence ``\textbf{Karim teaches a course}" is less than that of ``\textbf{a course Karim teaches}".\\
	\Square\ Given a bigram model trained on standard English, the perplexity of the sentence ``\textbf{Karim teaches a course}" is higher than that of ``\textbf{a course Karim teaches}". \\
	\Square\ Given a bigram model trained on standard English, the perplexity of the sentence ``\textbf{Karim teaches a course}" is equal to that of ``\textbf{a course Karim teaches}".\\
	\Square\ Given a bigram model trained on standard English, the perplexity of the sentence ``\textbf{Karim teaches a course}" is less than that of ``\textbf{a course Karim teaches}".\\
	\hline
\end{longtable}

\subsection*{1.2.}

Given these training sentences: 

\begin{tabular}{|lllll|}
	\hline
	The courses are interesting && The course is interesting && He teaches a course \\
	He teaches some interesting courses && His course is interesting &&\\
	\hline
\end{tabular}

\begin{enumerate}
	\item using a bigram model with Laplace smoothing, calculate these expressions' probabilities: ``The courses are interesting", ``The courses is interesting" and ``The courses his interesting".
	\item What do you notice?
	\item Now, apply lemmatization on these sentences.
	Calculate these two expressions' probabilities:  ``The courses are interesting" and ``The courses is interesting".
	\item Is lemmatization useful for spell checking task using language models? Justify.
\end{enumerate}


\section*{2. PoS tagging}

\subsection*{2.1.}

Select right propositions:

\begin{longtable}{|p{.95\textwidth}|}
	\hline 
	\Square\ In PoS using HMM, transitions between states (PoS) are represented as an LM. \\
	\Square\ PoS using HMM can easily integrate heterogeneous features. \\	
	\Square\ MEMM can use those features used by HMM. \\
	\Square\ MEMM is a discriminative model like HMM. \\
	
	\hline
\end{longtable}

\subsection*{2.2.}

This is a hidden Markov model for PoS tagging: 

\begin{tikzpicture}[
	> = stealth, % arrow head style
	shorten > = 1pt, % don't touch arrow head to node
	auto,
	node distance = 3cm, % distance between nodes
	semithick, % line style
	font=\small
	]
	
	\tikzstyle{every state}=[
	draw = black,
	thick,
	fill = white,
	minimum size = 4mm
	]
	
	\node[circle,draw] (qN) {N};
	\node[align=left,draw] (qNe) [left of=qN] {P("fish"|N)=0.05\\P("river"|N)=0.03};
	\node[circle,draw] (qV) [right of=qN] {V};
	\node[align=left,draw] (qVe) [right of=qV] {P("fish"|V)=0.01\\P("eat"|V)=0.02};
	\node[circle,draw] (qD) [below of=qN] {D};
	\node[align=left,draw] (qDe) [left of=qD] {P("the"|D)=0.3\\P("a"|D)=0.2\\P("an"|D)=0.15};
	\node[circle,draw] (qP) [right of=qD] {P};
	\node[align=left,draw] (qPe) [right of=qP] {P("I"|P)=0.2\\P("he"|P)=0.1\\P("she"|P)=0.1};
	\node[] () [below of=qD, yshift=1.5cm] {$\pi(P, D, V, N) = (0.4, 0.3, 0.1, 0.2)$};
	
	
	\path[->] 	
	(qN) 	edge [loop above] node {0.3} ()
	edge [bend left] node {0.7} (qV)
	(qV) 	edge [loop above] node {0.1} ()
	edge [bend left] node {0.5} (qN)
	edge [bend left] node {0.4} (qD)
	(qD)	edge [bend left] node {1.0} (qN)
	(qP)	edge [bend right] node {1.0} (qV);
	
	\path[dashed] 	
	(qN) 	edge [] node {} (qNe)
	(qV) 	edge [] node {} (qVe)
	(qD) 	edge [] node {} (qDe)
	(qP) 	edge [] node {} (qPe);
	
\end{tikzpicture}

\begin{enumerate}
	\item Calculate these two probabilities: P(V D N | ``fish a fish") et P(N D N | "fish a fish")
	\item Given C(N) = 200, C(V) = 100 et C(D)=100, calculate these probabilities using Laplace smoothing: P(D|V), P(D|N) et P(N|D).
	\item Recalculate the probabilities of the first question using this smoothing.
	\item How much labeled expressions ``\textbf{V D N}" exist in our training dataset? The determiner ``\textbf{D}" does not happen at the end of a sentence. 
	
\end{enumerate}

\end{document}
